\chapter{Kỷ luật nghe khó chịu nhưng thiếu là toang}
\markboth{Kỷ luật nghe khó chịu nhưng thiếu là toang}{Kỷ luật nghe khó chịu nhưng thiếu là toang}

\section{Đi trễ vài phút mà làm như trọng án}
\begin{truyen}{Đi trễ vài phút mà làm như trọng án}{Classroom Talk}
\chuhoa{H}{ọc sinh: ``Em trễ có mấy phút thôi mà thầy làm như em phạm lỗi lớn lắm.''}
\textit{Student: ``I was only a few minutes late, but you are acting like I committed some major offense.''}\\[4pt]

\teacherpair{Một lần vài phút nghe nhỏ. Nhưng cộng lại thành thói quen coi nhẹ giờ giấc thì không nhỏ nữa.}{A few minutes once may sound small. But when those minutes add up into a habit of disrespecting time, it is no longer small.}

\studentpair{Nhưng em vẫn có mặt mà.}{But I still showed up.}

\teacherpair{Có mặt thân xác không có nghĩa là em tôn trọng nhịp làm việc chung.}{Being physically present does not mean you are respecting the shared rhythm of work.}
\end{truyen}

\begin{baihoc}
Kỷ luật nhiều khi không nằm ở con số phút, mà nằm ở thái độ coi thường hay coi trọng việc chung.
\textit{(Discipline is not only about minutes. It is about your attitude toward shared responsibility.)}
\end{baihoc}

\begin{ghinhoanh}
\textbf{Short English to remember:}

Small lateness becomes a pattern.

Patterns shape character.
\end{ghinhoanh}

\begin{tuhocnhanh}
1. Em trễ vì bất khả kháng hay vì quen tay coi nhẹ? \textit{(Are you late from real obstacles or casual carelessness?)}

2. Em cần sửa khâu nào để đến đúng giờ hơn? \textit{(What part of your routine must change to be on time?)}
\end{tuhocnhanh}
\ngancach

\section{Em không làm bài một bữa chứ có bỏ học đâu}
\begin{truyen}{Em không làm bài một bữa chứ có bỏ học đâu}{Classroom Talk}
\chuhoa{H}{ọc sinh: ``Em quên làm bài đúng một hôm thôi, đâu phải em bỏ luôn. Sao thầy làm căng vậy?''}
\textit{Student: ``I forgot to do the assignment for just one day. It is not like I dropped out of school. Why are you making such a big deal out of it?''}\\[4pt]

\teacherpair{Một hôm đó có phải là lần đầu em nới tay với chính mình không?}{Was that one day really the first time you went easy on yourself?}

\studentpair{Chắc... không.}{Probably... not.}

\teacherpair{Vậy điều đáng nói không phải một bài. Là cái thói quen tự tha dễ quá.}{Then the real issue is not one assignment. The real issue is the habit of forgiving yourself too easily.}
\end{truyen}

\begin{baihoc}
Nhiều sự xuống dốc bắt đầu bằng câu: ``một lần thôi mà.''
\textit{(Many declines begin with the phrase: just this once.)}
\end{baihoc}

\begin{ghinhoanh}
\textbf{Short English to remember:}

Just once can start a habit.

Easy excuses lower your standard.
\end{ghinhoanh}

\begin{tuhocnhanh}
1. Em đang bỏ sót việc vì mệt thật hay vì tự nới tay quá nhanh? \textit{(Are you slipping because of real fatigue or easy self-excusing?)}

2. Em cần nguyên tắc tối thiểu nào để không tụt? \textit{(What minimum rule would keep you from sliding backward?)}
\end{tuhocnhanh}
\ngancach

\section{Em nói nhỏ mà cũng bị nhắc}
\begin{truyen}{Em nói nhỏ mà cũng bị nhắc}{Classroom Talk}
\chuhoa{H}{ọc sinh: ``Em có làm ồn đâu thầy. Em chỉ nói nhỏ với bạn bên cạnh thôi mà.''}
\textit{Student: ``I was not being loud, sir. I was only talking softly to the person next to me.''}\\[4pt]

\teacherpair{Nếu năm người trong lớp đều nói nhỏ một chút, em nghĩ lớp còn nghe bài nổi không?}{If five people in this class each talk softly a little bit, do you think the class can still hear the lesson properly?}

\studentpair{Nghe hơi toang.}{That sounds like it would fall apart.}

\teacherpair{Đúng. Nhiều cái cá nhân thấy nhỏ, nhưng cộng lại phá rất mạnh.}{Exactly. Many things feel small on the individual level, but together they can be highly destructive.}
\end{truyen}

\begin{baihoc}
Không phải cứ nhỏ với riêng mình là nhỏ với tập thể.
\textit{(Something that feels small individually may not stay small collectively.)}
\end{baihoc}

\begin{ghinhoanh}
\textbf{Short English to remember:}

Small noise multiplied is big noise.

Shared space needs shared restraint.
\end{ghinhoanh}

\begin{tuhocnhanh}
1. Em hay chỉ nhìn tác động của mình hay nhìn tác động cộng dồn? \textit{(Do you look only at your own action or at the combined impact?)}

2. Em có đang dùng chữ ``chỉ'' để giảm nhẹ lỗi của mình không? \textit{(Do you use the word just to shrink your own mistake?)}
\end{tuhocnhanh}
\ngancach

\section{Không ai thấy thì lách một chút có sao}
\begin{truyen}{Không ai thấy thì lách một chút có sao}{Classroom Talk}
\chuhoa{H}{ọc sinh: ``Nếu không ai thấy thì lách luật chút cũng đâu hại ai. Quan trọng là đừng bị bắt thôi.''}
\textit{Student: ``If nobody sees it, then bending the rules a little does not really hurt anyone. The main thing is just not getting caught.''}\\[4pt]

\teacherpair{Câu đó nguy hiểm không phải vì nó ồn. Nó nguy vì nó cho thấy em đang lấy việc không bị phát hiện làm thước đo đúng sai.}{That sentence is dangerous not because it is loud, but because it shows that you are using not being caught as your standard for right and wrong.}

\studentpair{Thực tế ngoài đời cũng vậy mà thầy?}{But that is how real life works outside school too, right?}

\teacherpair{Ngoài đời có người sống vậy. Nhưng sống kiểu đó lâu ngày, người đầu tiên em làm hỏng là chính mình.}{There are people in real life who live that way. But if you live like that for too long, the first person you damage is yourself.}
\end{truyen}

\begin{baihoc}
Đúng sai không nên đo bằng chuyện có ai nhìn thấy hay không.
\textit{(Right and wrong should not be measured by whether anyone is watching.)}
\end{baihoc}

\begin{ghinhoanh}
\textbf{Short English to remember:}

Not getting caught is not the same as being right.

Character shows most in private choices.
\end{ghinhoanh}

\begin{tuhocnhanh}
1. Em có từng tự cho phép mình sai vì nghĩ không ai biết? \textit{(Have you ever allowed yourself a wrong action because no one would know?)}

2. Chuẩn mực nào em muốn giữ ngay cả khi không ai nhìn? \textit{(What standard do you want to keep even when no one is watching?)}
\end{tuhocnhanh}
\ngancach